\clearpage
\section{Discussion}
\label{sec:Discussion}

This thesis set out to support software architects in choosing between an LLM workflow and an agent for a single task in an agentic system. The results answer the research question in three parts. First, the choice matters, and practice lacks support for it. Second, a framework can support the choice in three steps. It first qualifies the task within its system context, then checks its control requirements and finally weighs the priorities of the architect, within hard limits that no weight can offset. It then attaches the conditions for building the chosen orchestration paradigm. Third, such support can work inside the architect's own LLM assistant, where the model interprets the description and code applies the fixed rules. Following \textcite{gregorPositioningPresentingDesign2013}, this chapter interprets these results against the objectives and relates them to the general problem of choosing an orchestration paradigm.

The interviews support the research gap. All twelve participants knew the choice from their own projects, and ten made it by experience, intuition, client preference or platform defaults, without a formal method (Section~\ref{subsubsec:TADF/episode1}). No participant described a procedure that combines task characteristics, comparative evidence, constraints and priorities. The experiments show that the choice has consequences. Per task archetype, the success rates of the two orchestration paradigms differed by up to 26.7 percentage points, and for single task requirements by far more. In six of eight task archetypes, the agent needed 1.8 to 7.7 times the tokens per successful execution (Section~\ref{subsubsec:TADF/results}). On WorkBench, Oracle, which chose the agent only where the agent alone succeeded, would have reached 75.9 instead of 66.2 percent. Choosing in hindsight for each template would have reached 73.2 percent (Section~\ref{subsubsec:TADF/episode2}).

Measured against the design objectives (Section~\ref{subsec:TADF/objective-results}), the strengths and gaps follow one pattern. Where the evidence came from controlled comparisons and tests, it is firmer. The decision core rests on paired executions, and code applies the rules reproducibly. Where the evidence depends on use, it is only formative. Practitioners found the logic understandable and asked for an assistant form. However, no independent architect has yet decided with TADF, and on WorkBench routing by TADF added no success beyond Always LLM workflow. The hypothesis therefore holds only in part. In the window, TADF gives traceable recommendations that are the same for the same answers, and practitioners asked for the assistant form that it now has. Its benefit on new tasks, however, remains to be shown.

Three choices of the research design proved valuable. Implementing both orchestration paradigms for the same task instances allowed a paired comparison at the task level, which earlier benchmarks did not offer. Asking practitioners before the final design turned the spreadsheet into a skill. Linking each rule to its evidence made it possible to retire rules whose evidence did not hold. The scope, in contrast, was broad: eight task archetypes, three models, deep-dive runs, two evaluation episodes and a working skill. This breadth mapped the decision but left several rules on thin evidence, and a narrower scope would have tested fewer rules more firmly. The approach still points in the right direction. The rules of thumb of practitioners matched its agent gates, and practice and research both call for explicit guidance (Section~\ref{subsec:discussion/practice}). In a young field, the broad map is also useful. It shows where the evidence is firm, where it is thin and where research can start. It also suggests how decision support for architects may look in the future: inside the tools that architects already use, and based on a versioned rulebook that grows with new evidence.

Section~\ref{subsec:discussion/findings} discusses the choice between the two orchestration paradigms (SQ1 and SQ2), Section~\ref{subsec:discussion/practice} decision support in practice (SQ3), Section~\ref{subsec:discussion/contribution} the design principles and Section~\ref{subsec:discussion/limitations} the limitations and future research.

\subsection{Choosing Between LLM Workflows and Agents}
\label{subsec:discussion/findings}

SQ1 and SQ2 asked how the two orchestration paradigms compare and which task characteristics and execution conditions are associated with their differences. Neither paradigm was better in general. That also underlines the current agentic archtecture discurse shown in \ref{sec:Background}.  The agent was more successful where the next step or the sources could not be fixed in advance, and the LLM workflow where written rules decided. In most task archetypes, the agent needed more tokens. Taken together, the results suggest that the choice is not one decision but a sequence of levels. At the level of the system context, the lifecycle frames the choice, for example one-off use or recurring production. At the level of the task boundary, G1 to G3 separate parts that need no LLM and split tasks that are too large. At the level of control, G4 and G5 decide whether only an agent can perform the task. In the weighted comparison, the priorities of the architect decide between the remaining options, within hard limits. At the level of implementation, the build form and its components decide whether the chosen orchestration paradigm succeeds. TADF follows this order, and the paradigm is only one layer of the architecture decision.

At present, the last two gates mark the clearest difference. G4 asks whether every next step can be fixed before the run, and G5 whether the sources are known in advance. Where later actions depended on earlier outcomes, the agent solved 13 of 15 executions and the LLM workflow none. Where a search result determined the next query, the agent solved 23 of 25 executions and the LLM workflow 14. Where the source of a needed fact was unknown, the result was 24 against 12 (Section~\ref{subsec:TADF/initial-artifact}). This is the strength of the agent: it chooses the next step from what it observes. Control also depends on what can be observed. Neither paradigm succeeded when a state changed silently, and the agent recovered in only 16 of 25 executions when the procedure required a new check (Section~\ref{subsubsec:TADF/results}). On WorkBench, capped searches hid records, which added the question about the planned reads to G4 (Section~\ref{subsubsec:TADF/episode2}). Practitioners used the same distinction in their rules of thumb (P02, P03, P06, P07). It sharpens practitioner guidance that reserves agents for tasks that need flexibility or where rule-based approaches fall short \parencite{schluntzBuildingEffectiveAI2024,openaiPracticalGuideAgents2026}. The number of actions alone, in contrast, did not decide. When a read set the number of actions, both paradigms solved all 25 executions, and the LLM workflow needed roughly a third of the agent's mean token cost. Where outcome branches were known in advance, a state-machine LLM workflow solved 19 and 14 of 25 executions, against 24 and 15 for the agent (Section~\ref{subsubsec:TADF/results}).

Outside these two gates, the choice is a trade-off, and here the deployment indicators mostly favor the LLM workflow. It applied written rules more reliably and returned the reference decision in all 75 executions, against 60 for the agent. It also needed fewer tokens per successful execution in six of eight task archetypes. This supports judging success together with cost \parencite{kapoorAIAgentsThat2024}. In retrieval and planning, the agent was more successful, but these differences remain uncertain. Under equal weights, the fixed capability scores favor the LLM workflow, so that only G4 or G5 can select an agent (Appendix~\ref{app:artifact/scoring}). The agent leads only on robustness to messy input and evolvability, and evolvability was derived rather than measured. On the comparative path, the priorities of the architect thus decide. TADF recommends the LLM workflow when cost, safe failure and schema guarantee weigh highly. It recommends the agent only when robustness to messy input and evolvability weigh well above them. The LLM workflow has clear weaknesses, too. Predefined control does not ensure correct execution. On WorkBench, the revised LLM workflow was more successful than the ReAct agent with domain tools but caused 119 side-effect errors against 59 (Section~\ref{subsubsec:TADF/episode2}). Its success also depends on the build, as two implementations of the LLM workflow reached 41.6 and 66.2 percent.

These findings are consistent with what practice reports at the system level. Cost and reliability limit the use of agentic systems (Section~\ref{subsec:intro/motivation}), and 68 percent of deployed agents run at most ten steps before a person intervenes \parencite{panMeasuringAgentsProduction2025}. Gartner expects that over 40 percent of agentic AI projects will be canceled by the end of 2027. It names rising costs, unclear business value and inadequate risk controls, and it notes that many use cases presented as agentic do not need an agentic implementation \parencite{gartnerGartnerPredictsOver2025}. P06 reported clients who asked for an agent because the term was popular, although a simpler solution would have sufficed. Model providers likewise advise starting with the simplest solution, because agents trade latency and cost for task performance \parencite{schluntzBuildingEffectiveAI2024,openaiPracticalGuideAgents2026}. Research on workflow optimization recommends starting from a constrained static structure and reserving changes during execution for settings in which execution reveals information that a plan cannot anticipate \parencite{yueStaticTemplatesDynamic2026}. More agents do not help by default either. Under a shared protocol, only one of six multi-agent systems exceeded a single-agent baseline, and its small lead was within the reported uncertainty \parencite{fuDoMoreAgentsHelp2026}. For workplace automation in production, the evidence therefore suggests a default for current models: an LLM workflow, unless G4 or G5 applies. Where these gates apply, practice tends to bound the agent or to keep a person in the loop. Practitioners kept agents to bounded parts of a workflow (P02--P04, P09), and P10 turned validated agent solutions into workflows. Research follows the same path: \textcite{abuzakukOptimizingAgenticWorkflows2026} compile recurring tool-call sequences of an agent into deterministic composite tools. In practice, the agent is thus less a rival of the LLM workflow than a bounded component for the parts that cannot be fixed in advance.

Before production, the agent has a clear strength. In the seven task archetypes with tools, the same prebuilt ReAct agent reached 70.7 to 100 percent task success. Each LLM workflow, in contrast, had to be built as a graph for its individual task archetype (Sections~\ref{subsubsec:TADF/paradigms} and~\ref{subsubsec:TADF/results}). The agent needs no procedure designed for the task. For non-recurring tasks, proofs of concept and daily work in which the goal is not yet clear, this saves design effort that an LLM workflow repays only through repeated use. Practitioners distinguished one-off from recurring use and prototypes from production (P04, P09, P12), and P10 turned validated agent solutions into workflows. The agent can thus explore the path, and the LLM workflow can fix it once the task recurs. The build effort itself was not measured (Section~\ref{subsec:discussion/limitations}).

This boundary is not fixed. With GPT-5.2 instead of GPT-5.4-nano, the agent's gap to the LLM workflow in compliance narrowed from 23.3 to 6.7 percentage points (Section~\ref{subsubsec:TADF/results}). On a revised WorkBench, the best agent reached 97.7 percent, although that study also changed the execution interface and corrected benchmark defects \parencite{stylesWorkBenchRevisitedWorkplace2026}. In simulated dialogues, a model that executed a procedure from its prompt reached higher judged task success than a code-controlled graph, but at higher token cost \parencite{dennisInContextPromptingObsoletes2026}. The gains of multi-agent designs also vary with task structure and model capability \parencite{kimScienceScalingAgent2025}. If stronger models close the success gap, cost, traceability and safeguards would weigh more. Falling prices could weaken the cost argument, but token use for complex tasks has so far grown faster than the price per token has fallen \parencite{tinkoffStateAI2026}. The default above is therefore a statement about current models, not a lasting rule.

The breadth of the scope also limits how far these conclusions reach. Each task archetype had 15 task instances, and the deep-dive runs examined only action execution. Part of each difference may stem from the implementations rather than from the paradigm, because the paired systems also differed slightly in prompts and deterministic components (Appendix~\ref{app:protocol/implementation}). G4, G5 and the cost difference rest on clear and repeated differences. Other rules rest on single comparisons, and two task-fit criteria were retired when their evidence did not hold (Section~\ref{subsubsec:TADF/artifact-correction}). The findings therefore support a direction rather than a settled rule set: predefined control first, and model control where the next step or the sources cannot be fixed in advance.

\subsection{Architecture Decision Support in Practice}
\label{subsec:discussion/practice}

SQ3 asked how relevant the distinction is in practice and how far TADF supports architects. The distinction is relevant, and both the interviews and the literature call for support of this kind. All participants knew the choice from their projects, and P01, P02 and P06 saw value in making it explicit. Half of the participants already plan or build automations with LLM assistants or coding agents (P04--P07, P11, P12). Five supported an LLM assistant or skill that they could use directly in their work (P05, P06, P08, P11, P12). Engineering teams often lack clear guidance on how to put agentic systems into operation \parencite{bandaraPracticalGuideAgentic2026}. A proposal for AI-based architecture decision support likewise combines explained recommendations with decision records, but reports only a preliminary evaluation \parencite{arifAIDrivenDecisionSupport2025}. TADF shows one way to meet this demand. The consultation starts from the architect's own description, the assistant proposes answers with reasons, code applies the fixed rules and the architect confirms. This allocation follows the assisting archetype of agentic information systems, in which the human keeps the main responsibility \parencite{holldackAgenticInformationSystems2026}. As LLM assistants become common tools of architects, such support may well spread. Its effect on the decisions of independent architects, however, is not yet tested.

The interviews also explain why the framework grew large. Performance alone did not settle the choice. Approved platforms and data protection could rule out a design (P07, P11), and data readiness could constrain it (P12). Approval before irreversible actions was a hard limit, not a preference (P02, P03, P10, P12). Such safeguards are not always built: among 6,003 public n8n templates, only 167 contained a human-mediated path before external actions \parencite{tangCharacterizingLargeLanguage2026}. Most of these factors entered TADF. The system context frames the consultation, mandatory markings and the approval rule act as hard limits, and environment checks follow each route. Data readiness remains for the architect to check. More coverage means more effort, a tension known from architecture decision methods \parencite{falessiDecisionmakingTechniquesSoftware2011} that the window eases. Even so, the current evidence does not allow a certain decision for every task. Further research will likely add factors, such as multi-agent coordination, other model families or security requirements. Design science research with generative AI faces a rapid erosion of prescriptive knowledge \parencite{zurheidenDesignScienceResearch2026}. TADF is therefore best understood as a living framework. In medicine, living guidelines update a single recommendation as soon as new evidence becomes available, so that the unit of update is the recommendation, not the whole guideline \parencite{aklLivingSystematicReviews2017}. TADF was developed in this way. Each rule names its evidence, each evaluated state was frozen, and the development log records every change, including the retirement of two task-fit criteria. Continuing this practice beyond the thesis would keep the framework current.

The task level is central to this support, but it is hard to delimit in practice. A task is the unit that receives a route, yet the same work can form one task or several. In the text and the window consultation of the CrewAI email scenario, draft writing was split into two tasks once and kept as one task once. The difference came from the architect's gate answers (Appendices~\ref{app:transcripts/audit} and~\ref{app:transcripts/window}). Participants found the unit of assessment unclear (P02, P03), and P11 doubted that LLM workflows and agents can always be separated cleanly for one task. TADF limits the problem with explicit split rules and a limit of two split-or-simplify operations per task lineage, but it cannot define the tasks for the architect. Making task boundaries tangible, for example through worked examples or links to existing process models, is a practical need for future versions.

Much recent research addresses the choice at run time. Routers select a model, a reasoning strategy or a single-agent or multi-agent configuration for each incoming request \parencite{gaoSingleagentMultiagentSystems2025,suDifficultyAwareAgenticOrchestration2025,zhouSelectthenSolveParadigmRouting2026}. Search methods generate and test agent designs automatically \parencite{huAutomatedDesignAgentic2024}. These approaches react to each request or need executions of candidate designs. TADF acts earlier, when the architect decides what to build and must justify the choice against constraints and priorities. The two forms can complement each other. Design-time advice could fix the acceptable options, and a validated router could choose among them at run time. Decision support for architects may thus develop into a living, assistant-based consultation that combines evidence-based rules at design time with selection at run time and learns from the decision records of deployed systems. This combination has not been evaluated but is seen as a good way for future decisions support systems.

The demonstration also showed where the assistant must not decide alone. In the walkthrough of 27 August 2026, the assistant changed score contributions and omitted a required step despite fixed documents (Appendix~\ref{app:transcripts/audit}). The final skill therefore lets code apply the rules. Tests confirmed that this code reproduces every tested route, and in test consultations the window cut the consultation time by about 75 to 80 percent (Section~\ref{subsubsec:TADF/artifact-correction}). TADF thus applies the main finding of this thesis to itself: the model interprets, and predefined code controls the procedure. The tests cover the rule code, not the interpretation of descriptions, and where a host cannot show interactive pages, the text mode again leaves rule application to the model. Human confirmation remains a weak point. Professionals may accept or reject AI advice without examining its basis \parencite{jussupowAugmentingMedicalDiagnosis2021,lebovitzEngageNotEngage2022}, so a confirmed proposal is not necessarily a checked one.

For practice, three implications follow. Software architects can treat the LLM workflow as the default for current models, use an agent where G4 or G5 applies and implement the mandatory minimum. Clients and managers can ask whether a task needs an agent before they pay for one (P06). Builders of assistants and automation platforms can keep rule application in code and make approval paths easy to add.

\subsection{Design Principles}
\label{subsec:discussion/contribution}

Together with the comparative evidence, the decision method of TADF and four design principles form the knowledge contribution of the thesis. The catalogue of 18 agent design patterns of \textcite{liuAgentDesignPattern2025} links design choices to qualitative trade-offs. TADF addresses a narrower choice but connects it to paired executions. In the terms of \textcite{gregorPositioningPresentingDesign2013}, the work is an improvement: it offers an explicit, evidence-linked procedure for a known decision that practice addresses with heuristics. The skill is a situated instantiation at Level~1. The decision method and the principles are nascent design knowledge at Level~2. The method comprises gates before scoring, fixed capability scores and design knowledge with a mandatory minimum. A well-developed design theory is not claimed. The principles were derived from the design and the formative evaluation of TADF, which instantiates them. Whether other designers find them understandable and useful, and whether instances built from them improve decisions, remains to be tested, ideally in contrasting settings \parencite{australiannationaluniversityaustraliaResearchPerspectivesAnatomy2020,seidelDesignPrinciplesSensemaking2018}.

 The four principles describe decision support for the choice between LLM workflows and agents, of which TADF is one instance. They follow the schema of \textcite{australiannationaluniversityaustraliaResearchPerspectivesAnatomy2020}: for implementers to achieve or allow an aim for users in a context, employ mechanisms, involving enactors, because of a rationale. The implementers are researchers and developers who build such decision support. The users are software architects who decide before implementation, and the enactors are the architect, an LLM assistant and code. All four principles share one context, which also sets their boundary conditions: the task-level choice in agentic workplace automation with current LLMs.

DP1 to DP3 address, in this order, the three components of a DSS (Section~\ref{subsec:theory/decision-support}): the decision model, the data (here the evidence behind the rules) and the dialogue with the user. DP4 adds the division of work between the assistant, code and the architect, a question that arises when an LLM conducts the dialogue. The principles take up the design objectives of Section~\ref{subsec:TADF/objectives} (DP1: O1, O2, O7; DP2: O3, O6, O8; DP3: O7; DP4: O4), but not O5, because better decisions on new tasks remain unestablished. DP1 to DP3 use allow, because the architect and the assistant act in them, and DP4 uses achieve, because code turns the same input into the same output \parencite{australiannationaluniversityaustraliaResearchPerspectivesAnatomy2020}. Each principle gives its statement, mechanism, rationale, instantiation in TADF and limits.

\paragraph{DP1: Decide from coarse to fine.}
\emph{To allow software architects to reach a well-founded choice of the orchestration paradigm with bounded effort, decision support should resolve the choice in levels, from coarse cut-offs to fine weighting, within hard limits that no weight can offset.} The principle uses the first four levels of Section~\ref{subsec:discussion/findings} in their order: the system context with its lifecycle, the task boundary, the control requirements and the priorities. Coarse levels can end clear cases early, and finer levels refine the rest. Hard requirements, such as approval before consequential actions, hold on every route. In architecture decisions, constraints and minimum levels decide which alternatives are feasible before preferences rank them \parencite{falessiDecisionmakingTechniquesSoftware2011,amellerLinkingQualityAttributes2012}, and task-technology fit points to the control that a task needs \parencite{goodhueTaskTechnologyFitIndividual1995}. Empirically, the lifecycle mattered to practitioners (P04, P09, P12), G4 and G5 separated the paradigms most clearly, and the priorities decided the remaining cases (Section~\ref{subsec:discussion/findings}). TADF instantiates the task boundary, the control requirements and the priorities as G1 to G5 and a weighted comparison with mandatory markings. Its system context, however, shapes only the design knowledge and does not yet end a decision early. The order of the levels and its effect on effort were not tested against alternatives.

\paragraph{DP2: State the implementation behind every recommendation.}
\emph{To allow software architects to judge whether a recommendation fits their build, decision support should state the implementation, the model and the conditions for which its evidence holds, and attach what a build needs to match them.} DP2 covers the last level of Section~\ref{subsec:discussion/findings}, the implementation, because the labels LLM workflow and agent are not enough. Evidence holds only for the tasks, technologies and conditions it covers (Section~\ref{subsec:theory/decision-support}). On WorkBench, implementations under one label differed widely: four configurations of the benchmark's agent reached 43.0 to 60.9 percent, and two implementations of the LLM workflow 41.6 and 66.2 percent (Appendix~\ref{app:eval2/results}). Several results also changed with the model (Section~\ref{subsubsec:TADF/results}). Stating these conditions lets researchers compare and accumulate results, and it shows when a recommendation must be revised, as when TADF retired two task-fit criteria. TADF states that its evidence on agents comes from a ReAct agent and that the model can change the result. It names a leading build form with its build rules, and it asks for a commitment to the mandatory minimum. The effect of single components was not isolated, and environment checks follow the route instead of being asked beforehand.

\paragraph{DP3: Provide the consultation where architects already work.}
\emph{To allow software architects to use decision support in their daily work, decision support should provide its consultation within the environments they already use.} Currently, these include LLM assistants and coding agents, which can load decision support as a skill, a reusable package of instructions, reference documents and code. Users of decision support can often decide whether to use it, so ease of use matters, and explanations that require less effort are used more \parencite{spragueFrameworkDevelopmentDecision1980,gregorExplanationsIntelligentSystems1999}. Half of the participants already plan or build automations with LLM assistants or coding agents, and four did not want to fill in a separate spreadsheet for every decision (Section~\ref{subsubsec:TADF/episode1}). Skills are spreading as a format for these assistants: one marketplace listed 40,285 public skills, with a strong concentration on software engineering workflows \parencite{lingAgentSkillsDataDriven2026}. TADF instantiates the principle as a skill whose window cut the consultation time by about 75 to 80 percent in test consultations comared to naormal conversation questions (Section~\ref{subsubsec:TADF/artifact-correction}). Use by independent architects was not tested. The principle leaves the form open, and a skill requires assistants that can load it.

\paragraph{DP4: Let code apply the rules.}
\emph{To achieve recommendations for software architects that follow the rulebook and are the same for the same answers, decision support should let code apply its rules, while the assistant proposes the answers and the architect confirms them.} From documents alone, the assistant deviated from the rules in the walkthrough (Section~\ref{subsec:discussion/practice}), an instance of the fragile internal consistency of generative AI artifacts \parencite{zurheidenDesignScienceResearch2026}. Code, in contrast, applies the rules the same way every time. The experiments point the same way: written rules were applied more reliably under predefined control (Section~\ref{subsec:discussion/findings}). In TADF, the window applies the rules in code. Tests showed that it reproduces every route of the text consultations, agrees with an independent implementation of the rule text on 20,000 random answer sets and rules out the deviations of the walkthrough (Appendix~\ref{app:artifact/files}). In the window consultations, the architect kept 58 of the 64 proposed answers and changed six (Appendix~\ref{app:transcripts/window}). The limits of Section~\ref{subsec:discussion/practice} remain: the tests do not cover the interpretation of descriptions, the text mode leaves rule application to the assistant, and a confirmed proposal is not necessarily a checked one.

\subsection{Limitations and Future Research}
\label{subsec:discussion/limitations}
\label{subsec:discussion/future}
The results are limited in six areas. Two concern all results: the pace of the field and the independence of the research. Four concern single parts: the taxonomy, the task as the unit of decision, the size of the evaluation, and the benchmark mode together with the scoring model and the build effort. 

Two limitations apply to all results. First, the field moves fast. New models appear several times a year, and several results already depended on the model (Section~\ref{subsec:discussion/findings}). Every capability score and rule is therefore a snapshot. A fixed regression set that is rerun when models or implementations change would keep the evidence current. Second, development and evaluation were not fully independent. Tasks, implementations and rules were refined together, and one analyst built the taxonomy and condensed the interviews. The builder of TADF also took the architect role in the pilots, the demonstration and iterations. The results therefore show what the instrument can do in the hands of its builder. Independent replication is the most direct remedy.

The taxonomy covers a selected part of workplace automation. Thirteen of 48 benchmark sources were included, and the 120 task instances were constructed for this thesis (Appendix~\ref{app:taxonomy}). The deep-dive runs concern only action execution, the agent of the experiments was a ReAct loop, and the main experimental comparison used three OpenAI models. The pattern across task archetypes is therefore descriptive, not a ranking of the paradigms. Future research should code task requirements independently and add task families with open sources and long chains of outcome-dependent steps. It should also include other model families, agent designs and multi-agent coordination.

The task as the unit of decision is the most fragile construct. It makes the choice manageable, but its boundary depends on judgment, and a different boundary can change the route without any change in the work (Section~\ref{subsec:discussion/practice}). Research should study how architects delimit tasks in real systems and test whether independent architects draw the same boundaries for the same description. From such studies, criteria or reference examples for task boundaries could be derived.

The evaluation was small. Ten interviews and two written responses from a non-probability sample gave formative feedback. In the interviews, the interviewer explained the Decision Sheet before participants assessed it. Each task archetype had 15 task instances, and GPT-5.2 ran each task instance once per paradigm. One outcome check in task archetype F rejected even complete executions, although excluding it does not change the direction or the test result for F (Appendix~\ref{app:results/contrasts}). On WorkBench, one model was used, the development tasks remained in the test set, and the LLM workflow was revised after observed failures (Appendix~\ref{app:eval2/limits}). The lead of the LLM workflow over the strongest agent also remains uncertain, because the task instances of a template are not independent. An LLM judge with limited calibration and incomplete token records further limit some comparisons (Appendices~\ref{app:protocol/judges} and~\ref{app:protocol/metrics}). Future evaluations should freeze the framework and the implementations before testing, use independently selected task families as a hold-out set and repeat executions across model families.

Most importantly, no independent architect has yet used TADF to make a decision. In Episode~2, the benchmark mode bypassed the exits of G1 to G3 and used default priorities, under which only G4 or G5 could select an agent. WorkBench also favors predefined control, because all information lies in known databases, and it had informed the construction of earlier task instances. The episode therefore tested a narrow part of the procedure. It showed no routing value beyond Always LLM workflow and could not show whether TADF selects an agent where one is needed. TADF is designed for software architects, who bring constraints and priorities that a benchmark does not contain. The next study should therefore let independent architects decide comparable design cases with TADF, with their usual practice and with a simple decision aid. Assessors who do not know the decision method should judge the decisions and the success and cost of the resulting prototypes against criteria fixed in advance. Consultation time, the detection of misleading proposals and the stability across host models and paraphrased descriptions should be measured as well. In terms of FEDS, this would move the evaluation from formative and artificial settings to summative and naturalistic ones \parencite{venableFEDSFrameworkEvaluation2016}. Finally, the scoring model is a design proposal. Capability bands, fixed task-fit weights and aggregation rules were chosen for the instrument, not estimated, and their sensitivity to other thresholds was not tested (Appendix~\ref{app:artifact/scoring}). Sensitivity analyses would show how robust the recommendations are. The build effort of the two paradigms was not measured either. TADF keeps it out of the scoring and uses the system context, such as one-off or recurring use, only to scale the build effort in its design knowledge. Because coding agents now help to build LLM workflows (P04, P07), the gap in build effort is hard to assess and may shrink. Token costs also exclude development and maintenance. Future research should measure the build and maintenance effort of both paradigms, with and without coding assistants.

These limits do not remove the value of the work. In a young field, the thesis provides paired evidence at the task level, a working instrument and design principles that future research can test, extend and revise.